\documentclass[11pt,a4paper]{article}
\usepackage[margin=26mm]{geometry}
\usepackage[T1]{fontenc}
\usepackage{lmodern}
\usepackage{amsmath,amssymb,amsthm,mathtools,booktabs,microtype}
\usepackage[hidelinks]{hyperref}
\usepackage{enumitem}
\setlength{\parindent}{0pt}
\setlength{\parskip}{5pt plus 1pt minus 1pt}
\setlength{\emergencystretch}{2em}
\allowdisplaybreaks[1]
\newtheorem{theorem}{Theorem}[section]
\newtheorem{lemma}[theorem]{Lemma}
\newtheorem{proposition}[theorem]{Proposition}
\newtheorem{corollary}[theorem]{Corollary}
\theoremstyle{remark}
\newtheorem{remark}[theorem]{Remark}
\newcommand{\C}{\mathbb C}
\newcommand{\R}{\mathbb R}
\newcommand{\norm}[1]{\left\lVert#1\right\rVert}
\newcommand{\eps}{\varepsilon}
\newcommand{\om}{\omega}
\newcommand{\poly}{\mathcal P}
\newcommand{\defeq}{\mathrel{:=}}
\hypersetup{pdftitle={IE-16: A counterexample and the failure of every universal subset constant},pdfauthor={Sidney Holden},pdfsubject={Complete negative resolution with an exact nine-point certificate}}
\begin{document}
\begin{center}
{\LARGE\bfseries IE-16: A counterexample and the failure of every universal subset constant\par}
\par\vspace{8pt}
{\large Sidney Holden\par}
{\small Center for Computational Biology, Flatiron Institute\par}
{\small Simons Foundation\par}
{\small 12 September 2026\par}
\end{center}
\vspace{5pt}

\begin{abstract}
The proposed $4/\pi$ subset bound is false. Let $\om=e^{2\pi i/3}$ and
$L=\{\om^a+10^{-3}\om^b:0\le a,b\le2\}$. At degree $k=4$, this nine-point
set satisfies
\[
 \frac{M_4(L)}{\max_{S\subset L,\ |S|=5}M_4(S)}>\frac{13}{10}>\frac4\pi.
\]
The proof is analytic and covers every subset by its cluster occupancy.
An additional exact-arithmetic certificate checks all 126 subsets and proves
optimality of the full-set polynomial by positive weighted orthogonality.
The actual ratio is approximately $1.33299891979475829424$.
A three-cluster amplification lemma gives a stronger conclusion: no finite
constant, independent of dimension, degree and the points, can satisfy the
proposed type of bound. For every positive integer $r$ and every $\eta>0$,
examples with $n=3^r$ and $k=(3^r-1)/2$ have ratio greater than
$(2/\sqrt3)^r-\eta$.
\end{abstract}

\section{The question and the conclusions}
For a finite set $E\subset\C\setminus\{0\}$, put
\[
 M_k(E)=\min\{\norm{p}_E:p\in\C[z],\ \deg p\le k,\ p(0)=1\},
 \qquad \norm{p}_E=\max_{z\in E}|p(z)|.
\]
For $|E|\ge k+1$, define
\[
 B_k(E)=\max_{S\subseteq E,\ |S|=k+1}M_k(S),
 \qquad R_k(E)=\frac{M_k(E)}{B_k(E)}.
\]
All denominators used here are positive: a polynomial of degree at most $k$
with value one at zero cannot vanish on $k+1$ distinct nonzero points.
The finite-dimensional minima exist because evaluation on at least $k+1$
distinct points is injective.

IE-16 asks whether $R_k(E)\le4/\pi$ for every set of $n\ge3$ distinct nonzero
complex numbers and every $1\le k\le n-2$.\footnote{Canonical statement:
A. Townsend, \emph{Open Problems in Numerical Linear Algebra}, entry IE-16,
accessed 12 September 2026; Source 1 below.}
The original normal-GMRES paper formulates a universal-constant conjecture
and proposes $4/\pi$ as the candidate.\footnote{J. Liesen and P. Tich\'y,
\emph{The worst-case GMRES for normal matrices}, BIT 44 (2004), 79--98,
\S3.2.2, especially (3.16), pp. 91--92; Source 2 below.}
The following two theorems give negative answers, respectively, to the stated
constant and to the possibility of any universal finite replacement.

\begin{theorem}[Explicit counterexample]\label{thm:finite}
Let
\[
 \om=-\frac12+\frac{\sqrt3}{2}i,\qquad \eps=\frac1{1000},\qquad
 L_\eps=\{\om^a+\eps\om^b: a,b\in\{0,1,2\}\}.
\]
Then $L_\eps$ consists of nine distinct nonzero points, $k=4$ is admissible,
and
\begin{align}
 M_4(L_\eps)
 &=\frac{3003003000}{1001003001001}>\frac{299}{100000},\label{eq:M-exact-number}\\
 B_4(L_\eps)&<\frac{23}{10000}.\label{eq:B-bound-number}
\end{align}
Consequently $R_4(L_\eps)>13/10>4/\pi$.
\end{theorem}

\begin{theorem}[No universal finite constant]\label{thm:unbounded}
For every integer $r\ge1$ and every $\eta>0$, there is an admissible set $E$
of $n=3^r$ distinct nonzero complex numbers such that, at
$k=(3^r-1)/2$,
\[
 R_k(E)>\left(\frac2{\sqrt3}\right)^r-\eta.
\]
In particular, $\sup_{n,k,E}R_k(E)=\infty$, where the supremum is taken over
the assumptions and quantifiers of IE-16.
\end{theorem}

Theorem~\ref{thm:finite} alone completely resolves IE-16. Theorem~\ref{thm:unbounded}
is stronger; its proof does not rely on numerical experiments.
No minimality of nine as a counterexample size, nor an optimal growth rate in
$n$, is asserted.

\section{Two elementary interpolation formulas}
We will also use the monic approximation problem
\[
 T_d(K)=\min\{\norm{q}_K:q\in\C[z]\text{ is monic of degree }d\}
\]
for a finite set $K\subset\C$ with $|K|\ge d+1$. Here $K$ is allowed to contain
zero. Set
\[
 b_d(K)=\max_{U\subseteq K,\ |U|=d+1}T_d(U),
 \qquad \Gamma_d(K)=\frac{T_d(K)}{b_d(K)}.
\]
These quantities are positive and their minima exist by the same
finite-dimensional argument.

\begin{lemma}[Interpolation identities]\label{lem:interpolation}
For $S=\{z_0,\ldots,z_k\}\subset\C\setminus\{0\}$ with distinct points,
\begin{equation}\label{eq:lagrange-residual}
 M_k(S)=\left(\sum_{j=0}^k
 \prod_{h\ne j}\frac{|z_h|}{|z_j-z_h|}\right)^{-1}.
\end{equation}
For $U\subset\C$ with $|U|=d+1$, let
\[
 A_d(U)=\sum_{x\in U}\frac1{\prod_{y\in U\setminus\{x\}}|x-y|}.
\]
Then
\begin{equation}\label{eq:lagrange-monic}
 T_d(U)=\frac1{A_d(U)}.
\end{equation}
\end{lemma}

\begin{proof}
Write $\ell_j(z)=\prod_{h\ne j}(z-z_h)/(z_j-z_h)$.
Interpolation gives $1=p(0)=\sum_j\ell_j(0)p(z_j)$, so
$\norm{p}_S\ge1/\sum_j|\ell_j(0)|$.
Equality is attained by interpolating the values
\[
 p(z_j)=\frac{\overline{\ell_j(0)}}{|\ell_j(0)|\sum_h|\ell_h(0)|}.
\]
All $\ell_j(0)$ are nonzero. This proves~\eqref{eq:lagrange-residual}.
For the monic problem, replace evaluation at zero by extraction of the leading
coefficient. The leading coefficient of the cardinal polynomial at $x\in U$
is $\alpha_x=1/\prod_{y\ne x}(x-y)$; interpolation of
$\overline{\alpha_x}/(|\alpha_x|\sum_y|\alpha_y|)$ has leading coefficient
one and norm $1/\sum_y|\alpha_y|$. The matching lower bound follows from the
triangle inequality.
\end{proof}

\section{The exact nine-point counterexample}
For this section let $0<\eps\le1/10$, and retain the definition of $L_\eps$
in Theorem~\ref{thm:finite}. Its three clusters are
\[
 C_a=\{\om^a+\eps\om^b:0\le b\le2\},\qquad a=0,1,2.
\]
Points in one cluster have distance $\sqrt3\eps$. Points in different clusters
have distance at least $\sqrt3-2\eps>0$. Also
$1-\eps\le|z|\le1+\eps$, so the nine points are distinct and nonzero.

\subsection{The full-set minimum}
\begin{proposition}\label{prop:exact-M}
For $0<\eps\le1/10$, put
\[
 D=1+\eps+3\eps^2+\eps^3+\eps^4.
\]
Then
\begin{equation}\label{eq:full-M}
 M_4(L_\eps)=\frac{3\eps(1+\eps+\eps^2)}{D},
 \qquad
 p_*(z)=1-\frac{1+\eps}{D}z^3
\end{equation}
is a minimizing polynomial. Its modulus equals $M_4(L_\eps)$ at every point
of $L_\eps$.
\end{proposition}

\begin{proof}
The set is invariant under multiplication by $\om$. Averaging a feasible
polynomial over this group does not increase its norm, and leaves a polynomial
$1+cz^3$. Conjugation invariance then allows $c$ to be real, again by averaging.
The three distinct values of $z^3$ on $L_\eps$ are
\[
 W_0=(1+\eps)^3=U,\qquad W_1=V+iW,\qquad W_2=V-iW,
\]
where
\[
 V=1-\tfrac32\eps-\tfrac32\eps^2+\eps^3,
 \qquad W=\tfrac{3\sqrt3}{2}(\eps-\eps^2).
\]
Thus the squared minimax problem reduces to minimizing
\[
 \max\{g_0(c),g_1(c)\},\qquad
 g_0(c)=(1+Uc)^2,\quad g_1(c)=(1+Vc)^2+W^2c^2.
\]
Define
\[
 c_*=-\frac{1+\eps}{D},\qquad
 m=\frac{3\eps(1+\eps+\eps^2)}{D},\qquad
 H=1-3\eps+\eps^2-3\eps^3+\eps^4.
\]
Direct expansion gives
\begin{align*}
 1+Uc_*&=-m,\qquad g_0(c_*)=g_1(c_*)=m^2,\\
 g_0'(c_*)&=-2Um<0,\\
 \tfrac12 g_1'(c_*)&=V+c_*(V^2+W^2)
                 =\frac{3\eps(1+\eps)H}{2D}>0.
\end{align*}
For the last inequality,
$H\ge1-3\eps-3\eps^3\ge697/1000>0$.
If $c<c_*$, then $1+Uc<-m$, so $g_0(c)>m^2$.
If $c>c_*$, the convex quadratic $g_1$ is strictly increasing from $c_*$,
so $g_1(c)>m^2$. The global minimum is therefore $m^2$, attained at $c_*$.
All three cubic-image residual moduli are $m$, proving the final assertion.
\end{proof}

\subsection{Every five-point subset}
\begin{proposition}\label{prop:subset-upper}
For $0<\eps\le1/10$ and every five-point subset $S\subset L_\eps$,
\begin{equation}\label{eq:uniform-upper}
 M_4(S)\le
 U_\eps\defeq\frac{\sqrt3\eps(\sqrt3+2\eps)^3}{4(1-\eps)^4}.
\end{equation}
\end{proposition}

\begin{proof}
Let $s_a=|S\cap C_a|$. Up to permutation, the only possible profiles are
$(2,2,1)$, $(3,1,1)$ and $(3,2,0)$.
For any $z_j\in S$, the numerator of its absolute Lagrange coefficient
in~\eqref{eq:lagrange-residual} is at least $(1-\eps)^4$.

For profile $(2,2,1)$, each of the four points in the two two-point clusters
has one within-cluster denominator factor $\sqrt3\eps$ and three
cross-cluster factors at most $\sqrt3+2\eps$. Their four contributions give
\[
 \frac1{M_4(S)}\ge
 \frac{4(1-\eps)^4}{\sqrt3\eps(\sqrt3+2\eps)^3},
\]
which is~\eqref{eq:uniform-upper}.

For either profile containing a three-point cluster, the three points in
that cluster alone give
\[
 \frac1{M_4(S)}\ge
 \frac{3(1-\eps)^4}{(\sqrt3\eps)^2(\sqrt3+2\eps)^2}
 =\frac{(1-\eps)^4}{\eps^2(\sqrt3+2\eps)^2}.
\]
Thus
\[
 M_4(S)\le\frac{\eps^2(\sqrt3+2\eps)^2}{(1-\eps)^4}\le U_\eps,
\]
since $4\eps\le\sqrt3(\sqrt3+2\eps)$; indeed $4\eps\le2/5<3$.
These cases exhaust all subsets.
\end{proof}

\begin{proof}[Proof of Theorem~\ref{thm:finite}]
Substitute $\eps=1/1000$ in~\eqref{eq:full-M} to obtain the exact rational
value~\eqref{eq:M-exact-number}. It exceeds $299/100000$.
For the subset bound, expand
\[
 \sqrt3(\sqrt3+2\eps)^3
 =9+18\sqrt3\eps+36\eps^2+8\sqrt3\eps^3
 <9+36\eps+36\eps^2+16\eps^3.
\]
At $\eps=1/1000$, an entirely rational comparison gives
\[
 \frac{9+36/1000+36/10^6+16/10^9}{4(999/1000)^4}<\frac{23}{10}.
\]
Proposition~\ref{prop:subset-upper} now implies~\eqref{eq:B-bound-number}.
Finally,
\[
 R_4(L_\eps)>
 \frac{299/100000}{23/10000}=\frac{13}{10}>\frac4\pi,
\]
where the last inequality follows already from $\pi>31/10$.
Here $n=9$ and $1\le4\le7=n-2$, so every assumption in IE-16 is met.
\end{proof}

\subsection{The limiting ratio is \texorpdfstring{$4/3$}{4/3}}
\begin{proposition}\label{prop:four-thirds}
For the preceding family, $R_4(L_\eps)\longrightarrow4/3$ as $\eps\downarrow0$.
\end{proposition}
\begin{proof}
Equation~\eqref{eq:full-M} gives $M_4(L_\eps)/\eps\to3$.
For any fixed subset pattern of type $(2,2,1)$, each of the four points with
one within-cluster companion has absolute Lagrange coefficient asymptotic to
\[
 \frac1{(\sqrt3\eps)(\sqrt3)^3}=\frac1{9\eps}.
\]
The remaining coefficient is bounded. Consequently
$M_4(S)/\eps\to9/4$.
For the other two profiles, Proposition~\ref{prop:subset-upper}'s proof gives
$M_4(S)=O(\eps^2)$. There are only 126 fixed subset patterns, so the maximum
of their rescaled values tends to $9/4$. Taking the ratio proves the claim.
\end{proof}

\section{A three-cluster amplification lemma}
The distinction between residual and monic approximation is important in the
iteration below. We prove both limits rather than identify the two problems
without justification.

\begin{lemma}[Cluster amplification]\label{lem:amplification}
Let $d\ge1$, and let $K\subset\C$ be a finite set of $m\ge d+1$ distinct points.
Set $a_j=\om^j$, $j=0,1,2$, and for $\eps>0$ define
\[
 E_\eps=\bigcup_{j=0}^2(a_j+\eps K),\qquad k=3d+1.
\]
For sufficiently small $\eps$, these are $3m$ distinct nonzero points. As
$\eps\downarrow0$,
\begin{align}
 \frac{T_k(E_\eps)}{\eps^d}&\longrightarrow3^dT_d(K),&
 \frac{M_k(E_\eps)}{\eps^d}&\longrightarrow3^dT_d(K),\label{eq:full-limits}\\
 \frac{b_k(E_\eps)}{\eps^d}&\longrightarrow
       3^d\frac{\sqrt3}{2}b_d(K),&
 \frac{B_k(E_\eps)}{\eps^d}&\longrightarrow
       3^d\frac{\sqrt3}{2}b_d(K).\label{eq:subset-limits}
\end{align}
In particular,
\begin{equation}\label{eq:amplified-ratios}
 \Gamma_k(E_\eps)\longrightarrow\frac2{\sqrt3}\Gamma_d(K),
 \qquad R_k(E_\eps)\longrightarrow\frac2{\sqrt3}\Gamma_d(K).
\end{equation}
\end{lemma}

\begin{proof}
Let $P(z)=z^3-1$ and $t=T_d(K)$. Cluster separation and avoidance of zero
follow for sufficiently small $\eps$ because $K$ is finite and the three
centers are distinct and nonzero. We prove the four limits in three steps.

\medskip
\textbf{Step 1: upper bounds for the two full-set minima.}
Choose a monic minimizing polynomial
\[
 Q(x)=x^d+\sum_{q=0}^{d-1}q_qx^q,\qquad \norm Q_K=t.
\]
For the monic problem take $h(z)=z$; for the residual problem take
$h(z)=(-1)^d$. In both cases $|h(a_j)|=1$.
Let $F_0=P^dh$ and
\[
 A_j=(P'(a_j))^dh(a_j)=(3a_j^2)^dh(a_j),\qquad |A_j|=3^d.
\]
There is a unique polynomial $H_\eps$ of degree at most $3d-1$ with Hermite
data
\begin{equation}\label{eq:hermite-correction}
 \frac{H_\eps^{(q)}(a_j)}{q!}=A_jq_q\eps^{d-q},
 \qquad 0\le q<d,\quad j=0,1,2.
\end{equation}
Existence and uniqueness follow because the homogeneous system would have a
polynomial of degree at most $3d-1$ divisible by $P^d$.
The inverse interpolation map is fixed, independent of $\eps$. All data
in~\eqref{eq:hermite-correction} are $O(\eps)$, so every coefficient of
$H_\eps$ is $O(\eps)$.

Put $F_\eps=F_0+H_\eps$. Taylor expansion, uniformly for the finite set
$x\in K$, gives
\begin{equation}\label{eq:local-upper}
 \eps^{-d}F_\eps(a_j+\eps x)\longrightarrow A_jQ(x).
\end{equation}
Indeed, the terms of orders below $d$ are prescribed by
\eqref{eq:hermite-correction}, the order-$d$ coefficient tends to $A_j$,
and all higher terms vanish after rescaling.
In the monic case $F_\eps$ is monic of degree $3d+1$.
In the residual case $F_\eps(0)=1+O(\eps)$, so division by $F_\eps(0)$ gives
a feasible residual polynomial without changing the limit.
Thus both limsups in~\eqref{eq:full-limits} are at most $3^dt$.

\medskip
\textbf{Step 2: matching full-set lower bounds.}
Let $f_\eps$ be a minimizer of either full-set problem. Step 1 gives
$\norm{f_\eps}_{E_\eps}=O(\eps^d)$.
We first justify boundedness of its coefficients, which is needed for the
confluent limit.

Choose distinct $\xi_0,\ldots,\xi_d\in K$. At centers $a_0,a_1$ consider
divided differences of orders $q=0,\ldots,d$ on
$a_j+\eps\xi_0,\ldots,a_j+\eps\xi_q$; at center $a_2$ use orders
$q=0,\ldots,d-1$. There are $3d+2$ functionals, the dimension of polynomials
of degree at most $3d+1$. Their matrices on polynomial coefficients converge
to the Hermite matrix with multiplicities $(d+1,d+1,d)$.
That limiting matrix is invertible: a polynomial in its kernel would have
$3d+2$ zeros counted with multiplicity and degree at most $3d+1$.
The inverse matrices are consequently uniformly bounded for small $\eps$.

An order-$q$ divided difference of $f_\eps$ on these nodes has the form
\[
 \eps^{-q}\sum_{u=0}^q
 \frac{f_\eps(a_j+\eps\xi_u)}
 {\prod_{v\ne u}(\xi_u-\xi_v)},
\]
so it is $O(\eps^{d-q})$, hence bounded for $q\le d$.
Uniform boundedness of the inverse matrices proves boundedness of all
coefficients of $f_\eps$.

Now expand locally as
\[
 g_{j,\eps}(x)=\eps^{-d}f_\eps(a_j+\eps x).
\]
Terms of degree greater than $d$ tend to zero, because the unscaled
coefficients are bounded. The degree-at-most-$d$ truncations have bounded
values at $\xi_0,\ldots,\xi_d$, so their coefficients are bounded by fixed
Vandermonde interpolation. Passing to a subsequence gives both a coefficient
limit $f_0$ and degree-at-most-$d$ limits $g_j$.
For $q<d$ the local bounds imply
$f_\eps^{(q)}(a_j)=O(\eps^{d-q})\to0$. Therefore
\[
 f_0=P^dh,\qquad \deg h\le1,
\]
and the leading coefficient of $g_j$ is $(3a_j^2)^dh(a_j)$.
In the monic case $h$ is monic linear; in the residual case
$h(0)=(-1)^d$.

In either case $\max_j|h(a_j)|\ge1$. For $h(z)=z+c$ this follows from
\[
 \frac13\sum_{j=0}^2|a_j+c|^2=1+|c|^2;
\]
for $h(z)=(-1)^d+bz$ it follows from the analogous identity
$1+|b|^2$. By the definition of the monic minimum,
\[
 \norm{g_j}_K\ge3^d|h(a_j)|t,
\]
including the case of zero leading coefficient, when the right side is zero.
Taking the maximum and a subsequence attaining the original liminf gives a
lower bound $3^dt$ for both liminfs. This proves~\eqref{eq:full-limits}.

\medskip
\textbf{Step 3: all subset limits.}
Label a $(k+1)$-point subset by its fixed choices inside the three copies of
$K$, with occupancy $(s_0,s_1,s_2)$ and sum $3d+2$.
There are finitely many such labeled patterns.
If some $s_j\ge d+2$, a term of the monic Lagrange sum has order
$\eps^{-(s_j-1)}$ with a positive limiting coefficient.
The associated subset minimum is therefore $o(\eps^d)$.
The same holds for residual interpolation, because its additional numerator
$\prod_{h\ne j}|z_h|$ tends to one.

If all $s_j\le d+1$, their sum forces the profile $(d+1,d+1,d)$.
Let $U,V\subset K$ be the two $(d+1)$-point choices.
For a point in either full cluster, there are $d$ within-cluster denominator
factors and $2d+1$ cross-cluster factors. Since every distance between distinct
centers is $\sqrt3$, the rescaled monic Lagrange sum tends to
\begin{equation}\label{eq:subset-lagrange-limit}
 \frac{A_d(U)+A_d(V)}{(\sqrt3)^{2d+1}}
 =\frac{A_d(U)+A_d(V)}{3^d\sqrt3}.
\end{equation}
Contributions from the $d$-point cluster vanish after multiplication by
$\eps^d$. Formula~\eqref{eq:subset-lagrange-limit} is also the residual
limit, again because all numerator moduli tend to one.

The two full-cluster choices are independent. Their smallest possible sum
is $2\min_{|U|=d+1}A_d(U)=2/b_d(K)$, attained by choosing a minimizing subset
in each of those two clusters. Any $d$ points may be used in the third cluster.
Taking inverses in~\eqref{eq:subset-lagrange-limit}, and then the maximum over
the finitely many patterns, proves~\eqref{eq:subset-limits}.
Division yields~\eqref{eq:amplified-ratios}.
\end{proof}

\section{Iterating the amplification}
\begin{proof}[Proof of Theorem~\ref{thm:unbounded}]
Let $K_1=\{1,\om,\om^2\}$ and $d_1=1$.
Averaging gives $M_1(K_1)=1$, while
$\frac13\sum_j|a_j+c|^2=1+|c|^2$ gives $T_1(K_1)=1$.
Every two-point subset has monic and residual minimum $\sqrt3/2$, by
Lemma~\ref{lem:interpolation}. Hence
\[
 \Gamma_1(K_1)=R_1(K_1)=c,\qquad c=2/\sqrt3>1.
\]
We prove inductively that for every $r\ge1$ and every $\eta>0$ there is a set
$K_r$ of $3^r$ distinct nonzero points such that, for
$d_r=(3^r-1)/2$, both
\[
 \Gamma_{d_r}(K_r)>c^r-\eta,
 \qquad R_{d_r}(K_r)>c^r-\eta.
\]
The base case was just established. For the induction step, first choose
$K_{r-1}$ with monic ratio greater than
$c^{r-1}-\eta/(2c)$. Apply Lemma~\ref{lem:amplification} to this fixed set
and take the new scale $\eps$ sufficiently small that both limiting ratios
are approximated within $\eta/2$.
Their common limit is $c\Gamma_{d_{r-1}}(K_{r-1})>c^r-\eta/2$.
The new degree is
$3d_{r-1}+1=d_r$ and the new cardinality is $3\cdot3^{r-1}=3^r$.
This proves the induction.

The construction has the recursive form
\[
 K_r=\bigcup_{a=0}^2\bigl(\om^a+\eps_r K_{r-1}\bigr),
\]
where each scale is chosen after the inner set is fixed. No interchange of
an infinite-depth limit and a small-scale limit is used.
For every fixed $r$, all scales are positive and the resulting set is finite.
The scales may also be chosen rational, since the required strict inequalities
hold for every sufficiently small positive scale.
Finally $n=3^r\ge3$ and $k=(n-1)/2$ satisfies $1\le k\le n-2$.
Since $c^r\to\infty$, the residual subset ratios are unbounded.
\end{proof}

Along this sequence of dimensions, the supremum of the ratio is at least
$n^{\alpha}$, where $\alpha=\log(2/\sqrt3)/\log3>0$.
This is a lower-growth construction, not an assertion that the exponent is
optimal. It also shows why merely increasing the proposed numerical constant
cannot restore a dimension-independent theorem.

\section{An actual normal-GMRES witness and exact verification}
\subsection{Positive weights certify optimality}
The finite counterexample can be checked without relying on a nonlinear
optimizer or on an equality between two separately defined GMRES bounds.
Keep $D$ from Proposition~\ref{prop:exact-M}, and set
\[
 H=1-3\eps+\eps^2-3\eps^3+\eps^4,
 \qquad Q=(1+\eps)^2(1+\eps+\eps^2),
\]
\[
 \mu_0=\frac{H}{3D},\qquad \mu_1=\mu_2=\frac{Q}{3D},
 \qquad \nu_{a,b}=\frac13\mu_{(b-a)\bmod3}.
\]
For $0<\eps\le1/10$, these weights are positive and sum to one; the latter
follows from $H+2Q=3D$.
Writing $z_{a,b}=\om^a+\eps\om^b$, direct algebra gives
\begin{equation}\label{eq:weighted-orthogonality}
 \sum_{a,b}\nu_{a,b}\overline{p_*(z_{a,b})}\,z_{a,b}^{\ell}=0,
 \qquad \ell=1,2,3,4.
\end{equation}
For $\ell\not\equiv0\pmod3$ the sum vanishes by rotation. For $\ell=3$,
conjugation cancels the imaginary parts, and the identity is the cancellation
of
\[
 \mu_0 U(1+c_*U)+(1-\mu_0)\bigl(V+c_*(V^2+W^2)\bigr)=0
\]
using the expressions in Proposition~\ref{prop:exact-M}.

Let $m=M_4(L_\eps)$. For any feasible $p$, the difference $p-p_*$ has no
constant term, so~\eqref{eq:weighted-orthogonality} implies the exact identity
\begin{equation}\label{eq:pythagoras}
 \sum_{a,b}\nu_{a,b}|p(z_{a,b})|^2
 =m^2+\sum_{a,b}\nu_{a,b}|p(z_{a,b})-p_*(z_{a,b})|^2\ge m^2.
\end{equation}
Together with $|p_*(z_{a,b})|=m$, this is a second proof of the full-set value.

Now take the nonsingular normal matrix
\[
 A=\operatorname{diag}(z_{0,0},z_{0,1},z_{0,2},z_{1,0},\ldots,z_{2,2})
\]
and the unit initial residual with entries $(r_0)_{a,b}=\sqrt{\nu_{a,b}}$.
The defining residual minimization for GMRES at step four is exactly the
square root of the left side of~\eqref{eq:pythagoras}, minimized over feasible
$p$. Its value is $m$. Conversely, $p_*$ gives a residual of norm $m$ for
\emph{every} unit initial residual, since its nine diagonal residual entries
all have modulus $m$. Thus $m$ is indeed the worst-case step-four norm for
this explicit normal matrix.

\subsection{What the computational certificate checks}
The program \texttt{code/verify.py} requires only Python's standard library.
For $\eps=1/1000$, it represents a node as $a+b\om$ with rational $a,b$ and
uses
\[
 |a+b\om|^2=a^2-ab+b^2.
\]
It checks positivity and normalization of the nine weights, all nine residual
moduli squared, and each of the four complex identities
\eqref{eq:weighted-orthogonality} exactly.
For each of the 126 subsets it forms each squared absolute Lagrange weight
as a rational number and encloses its square root by integer square roots.
Rational alternating-series bounds in Machin's identity enclose $\pi$.
Every strict inequality is compared as a rational inequality, without a
floating-point tolerance.

The generated certificate gives
\begin{align*}
 M_4(L_{1/1000})&=0.002999994003011984982048\ldots,\\
 B_4(L_{1/1000})&=0.002250559965550380002031\ldots,\\
 R_4(L_{1/1000})&=1.332998919794758294237942\ldots,\\
 M_4(L_{1/1000})-\frac4\pi B_4(L_{1/1000})
 &>0.0001344920570754357296.
\end{align*}
The first three decimal displays are descriptive; the exact rational interval
endpoints are stored in \texttt{certificates/all\_subset\_certificate.json}.
The profile counts are 81 for $(2,2,1)$, 27 for $(3,1,1)$ and 18 for $(3,2,0)$.
The analytic proof above does not depend on this enumeration. The verification
program certifies the finite witness; it is not a formal proof-assistant
verification of the general amplification theorem.

\section*{Sources}
\begin{enumerate}[label={[\arabic*]},leftmargin=2em]
\item Alex Townsend, \emph{Open Problems in Numerical Linear Algebra},
IE-16, ``A sharp subset bound for worst-case normal GMRES.''
Canonical statement accessed 12 September 2026. The page inspected lists
``Partially resolved'' and a literature-check date of 10 September 2026.
\url{https://github.com/ajt60gaibb/OpenProblemsInNLA/blob/main/linear-systems-and-elimination/IE-16/README.md}
\item J\"org Liesen and Petr Tich\'y, \emph{The worst-case GMRES for normal
matrices}, BIT Numerical Mathematics \textbf{44} (2004), 79--98.
The subset formula is (3.10); the original conjecture and candidate constant
are in \S3.2.2, (3.16), pp. 91--92. Author-hosted full text:
\url{https://page.math.tu-berlin.de/~liesen/Publicat/LieTic04.pdf}
\end{enumerate}

\end{document}
